\documentclass{article}
\usepackage{ctex}
\usepackage{amsmath}
\usepackage{tikz}
\usepackage{unicode-math}
\usetikzlibrary{arrows.meta, decorations.markings,patterns}
\begin{document}

2.2-8 如题图所示,一电容器两极板都是边长为a的正方形金属平板,两板不是严格平行,而是有一夹角$\theta$.证明：当 $\theta\ll d/a$ 时,略去边缘效应,它的电容为

$$C=\varepsilon_0 \frac{a^2}d\Big(1-\frac{a\theta}{2d}\Big).$$






\newpage

2.2-15 四个电容器的电容分别为$C_1,C_2,C_3$ 和$C_4$,连接如题图,分别求：

(1)$AB$间的电容

$(2)DE$间的电容

$(3)AE$间的电容







\newpage

2.2-25题图中的电容$C_1,C_2,C_3$都是已知的,电容$C$是可以调节的.问当$C$调节到$A$、$B$两点的电势相等时,$C$的值是多少？










\newpage

2.2-32 (1)一平行板电容器两极板的面积都是$S$,相距为$d$,电容为$C=\frac{\varepsilon_0S}{d}.$当在两板上加电压$U$时,略去边缘效应,两板间电场强度为$E=\frac{U}{d}.$其中一板所带电荷量为$Q=CU$ ,故它所受的力为

$$F=QE=CU(\frac Ud)=\frac{CU^2}{d}.$$

这个结果对不对？为什么？

(2) 用虚功原理证明：正确的公式应为

$$F=\frac12\frac{Q^2}{Cd}=\frac12\frac{CU^2}d.$$







\end{document}